\section{Conditional-event coherence and sharp propagation: proofs}
\label{sec:annex-p15-rationality-4-4}

These are explicit finite semantic reconstructions of the chapter's
selected formal claims. They establish their stated constituent and
logical-independence scope; the source's richer nested-conditional result
retains its cited boundary \cite{Pfeifer2021ProbabilityLogic}.

\begin{annexitem}{annex:p15.rat44.coherence}{thm:p15-rat44-coherence-method}{Finite separation and zero-weight recursion}
\begin{proof}
For unconditional events a portfolio has gain \(s\cdot(q-p)\) at
the constituent with vector \(q\). If \(p\) is a convex combination
of the constituent vectors, the weighted mean of each such gain is
zero, so its range contains zero. Conversely let \(p\) be outside
their compact convex hull, and choose the closest hull point \(y\)
in Euclidean norm. Put \(d=p-y\ne0\). For every hull point \(q\),
the segment from \(y\) to \(q\) stays in the hull. Minimality at
its initial endpoint gives \(d\cdot(q-y)\leq0\). Hence
\(d\cdot(q-p)\leq-\|d\|^2<0\) at every constituent. The stake
vector \(d\) gives a sure loss, contradicting coherence.
This proves the convex-hull and linear-weight characterizations.

For a conditional family let
\(g(\omega)=(\mathbf1_{H_i}(\mathbf1_{E_i}-p_i))_i\).
On the full union of this family's conditioning events, the same
separation proof characterizes zero in the convex hull by portfolio
ranges containing zero on that full union. Coherence on each
portfolio's possibly smaller active union implies this necessary
first-stage condition, but the condition alone need not suffice for
subfamilies whose conditioning events get weight zero. That is why
the feasibility system in the first step is followed by recursion.
Coherence for all subfamilies implies feasibility also for the
zero-weight subfamily selected at each subsequent stage; thus coherent
assessments cannot be rejected by the procedure.

At any feasible stage, for every conditioning event not forced to have
zero weight, choose a feasible solution giving it positive weight.
Averaging these finitely many solutions produces one feasible measure
positive on all such events simultaneously. At least one exists: a
normalized measure on the union of the remaining events cannot assign
zero to every event in that union. Record this measure as a stratum
and recurse on the events forced to weight zero. Each iteration removes
at least one assessment, proving termination in at most the family size.
If the procedure accepts, every \(H_i\) has a first stratum giving
it positive weight, and at that stratum
\[
 \mu(E_i\mathbin{\&}H_i)=p_i\mu(H_i).
\]
For any nonzero portfolio choose the earliest stratum positive on any
of its conditioning events. Every later-first event has zero weight
there; every earlier-first event is absent from this portfolio by
earliestness; and each first-at-this-stage bet has zero mean gain.
The portfolio therefore has mean zero under this stratum, with
positive mass on its active union. Its gain cannot be strictly positive
or strictly negative throughout that union. This proves coherence.

This argument also proves a useful construction principle: any finite
sequence of normalized measures covering every assessed conditioning
event defines coherent prices by its first-positive-stratum ratios.
Later measures may put weight on events covered earlier; they do not
change those already determined prices. Such strata are mathematical
witnesses of conditional coherence, not an assumption of positive
ordinary mass on every possible event.
\end{proof}
\end{annexitem}

\begin{annexitem}{annex:p15.rat44.laws}{thm:p15-rat44-basic-laws}{Complement, identity, and ordinary conditional ratios}
\begin{proof}
On \(A\), the gain for \(A\mid A\) is \(s(1-p)\); coherence
for arbitrary signed stakes forces \(p=1\). For \(\sim A\mid A\)
the gain is \(-sp\), forcing \(p=0\). For complementary
conditionals with prices \(q,r\), the unit-stake sum is the constant
\(1-q-r\) on \(A\). It can avoid a strict sign only if
\(q+r=1\). The convex-hull representation on unconditional
constituents gives normalization, nonnegativity, and finite additivity,
hence the ordinary event identities.

Include \(A,A\mathbin{\&}B\), and \(B\mid A\) in one assessed
family. The first feasible constituent measure in the conditional
method has the specified unconditional event probabilities and
satisfies
\(\mu(A\mathbin{\&}B)-p(B\mid A)\mu(A)=0\).
If \(p(A)>0\), division gives the stated quotient. If \(p(A)=0\),
both unconditional terms are zero and the equation alone imposes
no conditional value. The recursive test still applies on \(A\):
if \(A\) logically implies \(B\), its price must be one; if their
intersection is impossible, the price must be zero. Null probability
is therefore not permission to disregard logical constraints.
\end{proof}
\end{annexitem}

\begin{annexitem}{annex:p15.rat44.paradox}{thm:p15-rat44-paradox}{Explicit endpoint witnesses and material implication}
\begin{proof}
Let the four \(A,B\) constituents be possible. For any desired
\(x,z\in[0,1]\), let the first measure be supported on \(\sim A\),
with weights \(x,1-x\) on \(\sim A\mathbin{\&}B\) and
\(\sim A\mathbin{\&}\sim B\). Let the next measure be supported
on \(A\), with weights \(z,1-z\) on
\(A\mathbin{\&}B\) and \(A\mathbin{\&}\sim B\).
The first-positive-stratum construction gives \(p(B)=x\) and
\(p(B\mid A)=z\), including both endpoints. Thus no narrower
conditional interval is possible.

The event \(\sim A\lor B\) contains \(B\), so its probability
is at least \(x\) and at most one. Every \(r\in[x,1]\) is
attained by ordinary weights \(x,r-x,1-r\) on
\(A\mathbin{\&}B\), \(\sim A\mathbin{\&}\sim B\), and
\(A\mathbin{\&}\sim B\), respectively; the fourth weight is zero.
These sum to one and give material-conditional probability \(r\).

Under positive-probability quotient conditioning, probability-zero
\(B\) has zero intersection with \(A\), forcing conditional zero;
probability-one \(B\) has the same intersection probability as
\(A\), forcing conditional one. For \(0<x<1\) and any target
\(z\in[0,1]\), choose sufficiently small \(a>0\) such that
\(az\leq x\) and \(a(1-z)\leq1-x\). Use weights
\(az,a(1-z),x-az,1-x-a(1-z)\) on the four constituents.
They are nonnegative, give \(p(A)=a>0,p(B)=x\), and give the
conditional quotient \(z\). This establishes the interior and
endpoint contrast precisely.
\end{proof}
\end{annexitem}

\begin{annexitem}{annex:p15.rat44.connexive}{thm:p15-rat44-connexive}{The declared probability interpretation and classical reductions}
\begin{proof}
For a contingent \(A\), both \(A\) and \(\sim A\) are admissible
conditioning events. The basic identities give the four Aristotle
assessments. Narrow-scope negation of its contrary-consequent
conditional replaces that consequent by its complement, giving the
probability-one identity conditional. For Boethius, the consequent of
the first principle is accepted exactly when
\(p(\sim\sim B\mid A)>h\), the identical condition to its
antecedent \(p(B\mid A)>h\). The second uses the identical
condition for \(\sim B\). These are implications between
high-assessment predicates, not a new nested conditional semantics.
If both opposed conditionals were accepted with threshold greater
than one-half, their probabilities would sum to more than one,
contradicting the complement law. This proves the Abelard reading.

With material implication, ``if \(\sim A\) then \(A\)'' is
\(A\lor A=A\), so its negation is \(\sim A\); the other
Aristotle negation is \(\sim\sim A=A\). The first Boethius
formula becomes
\[
 \sim(\sim A\lor B)\lor\sim(\sim A\lor\sim B)
 =(A\mathbin{\&}\sim B)\lor(A\mathbin{\&}B)=A.
\]
The second gives the same two disjuncts in reverse order. The
conjunction of opposed material conditionals is
\((\sim A\lor B)\mathbin{\&}(\sim A\lor\sim B)=\sim A\),
so its negation is \(A\). These explicit reductions prove the
classical comparison without conflating negation conventions.
\end{proof}
\end{annexitem}

\begin{annexitem}{annex:p15.rat44.noninformation}{thm:p15-rat44-noninformation}{All three noninformative argument forms}
\begin{proof}
For strengthening, use a first measure supported on
\(B\mathbin{\&}\sim P\), with \(F\)-probability \(x\), and a
second supported on \(B\mathbin{\&}P\), with \(F\)-probability
the arbitrary desired \(z\). First-positive-stratum prices give
\(p(F\mid B)=x,p(F\mid(B\mathbin{\&}P))=z\), for every \(x,z\in[0,1]\).

For contraposition, when \(x<1,z<1\), give unnormalized weights
\(x/(1-x),1,z/(1-z),0\) to
\(A\mathbin{\&}B,A\mathbin{\&}\sim B,
  \sim A\mathbin{\&}\sim B,\sim A\mathbin{\&}B\), respectively,
then normalize.
Both conditioning events have positive probability and their ratios
are \(p(B\mid A)=x\) and \(p(\sim A\mid\sim B)=z\).
For \(x=1,z<1\), put the first measure on \(A\mathbin{\&}B\), and the
second on \(\sim B\) with \(\sim A\)-probability \(z\).
For \(z=1,x<1\), put the first measure on \(\sim A\mathbin{\&}\sim B\)
and the second on \(A\) with \(B\)-probability \(x\).
For \(x=z=1\), equal ordinary weights on \(A\mathbin{\&}B\) and
\(\sim A\mathbin{\&}\sim B\) suffice. The stratum principle proves coherence
in the boundary cases as well.

For transitivity use a first measure on \(\sim A\mathbin{\&}B\)
with \(C\)-probability \(y\). Use a second measure on \(A\),
giving \(B\)-probability \(x\) and \(C\)-probability arbitrary
\(z\); for example use their two-coordinate product weights
\(xz,x(1-z),(1-x)z,(1-x)(1-z)\) on that stratum.
The resulting prices are \(p(B\mid A)=x,p(C\mid B)=y,p(C\mid A)=z\).
The product is only a chosen witness for attainability, not an
independence hypothesis imposed on the inference. The first stratum
makes \(A\) null, preserving the freedom of its conditional price
even when \(x=y=1\). Every value \(z\) has now been witnessed
for each argument form, proving the full noninformative intervals.
\end{proof}
\end{annexitem}

\begin{annexitem}{annex:p15.rat44.propagation}{thm:p15-rat44-propagation}{Sharp AND and cumulative bounds without independence}
\begin{proof}
For any coherent conditional family whose conditions are all \(A\),
the gain-vector feasibility argument on \(A\) supplies normalized
constituent weights representing its prices. Write
\(b=p(B\mid A),c=p(C\mid A),r=p((B\mathbin{\&}C)\mid A)\).
The four conditional constituent weights are
\(r,b-r,c-r,1-b-c+r\); their nonnegativity is exactly
\(\max\{0,b+c-1\}\leq r\leq\min\{b,c\}\).
For interval premises the smallest feasible \(r\) is
\(\max\{0,x_L+y_L-1\}\), attained at \(b=x_L,c=y_L\).
The largest is \(\min\{x_U,y_U\}\), attained at
\(b=x_U,c=y_U\). The displayed four weights construct each
endpoint, and convex mixtures of the two endpoint distributions
preserve the premise intervals and cover every intermediate value.
Put ordinary positive mass on \(A\) to realize these as coherent
conditional prices.

For CUT include the assessments
\(B\mid A,C\mid(A\mathbin{\&}B),C\mid A\)
and use feasibility on their active union \(A\). A representing
measure there has \(\mu(B)=x\), \(\mu(C)=z\), and
\(\mu(B\mathbin{\&}C)=y\mu(B)=xy\). Here the events are restricted
to the local constituent space \(A\). Hence
\(z=xy+\mu(\sim B\mathbin{\&}C)\), with
\(0\leq\mu(\sim B\mathbin{\&}C)\leq1-x\), proving the bounds.
For their sharpness, choose within \(A\) the four weights
\[
             xy,\quad x(1-y),\quad(1-x)t,\quad(1-x)(1-t),
             \qquad t\in[0,1],
\]
on \(B\mathbin{\&}C,B\mathbin{\&}\sim C,
      \sim B\mathbin{\&}C,\sim B\mathbin{\&}\sim C\).
If \(x>0\), they
give all prices by positive conditioning and yield
\(z=xy+(1-x)t\). If \(x=0\), the first measure is supported on
\(A\mathbin{\&}\sim B\), with \(C\)-probability \(t\). Add a second
measure on \(A\mathbin{\&}B\) with \(C\)-probability \(y\). The stratum
principle supplies the same prices and every \(z=t\), verifying
the otherwise problematic zero-conditioning endpoint. At \(x=1\)
the interval collapses correctly to \(\{y\}\).

Substitute \(A=\top\) in this construction. Then \(B\mid A\)
is unconditional \(B\), \(C\mid(A\mathbin{\&}B)\) is \(C\mid B\), and
\(C\mid A\) is unconditional \(C\). This proves MP and its
endpoint constructions, rather than assuming a ratio on null \(B\).
\end{proof}
\end{annexitem}
